\paragraph{Bilan signé et somme avec quotient entrant.}
\label{ampl-alfa-z-s-aclaracion}
La division positionnelle comporte deux étapes consécutives. Avec
\(K_j^{\mathrm{per}}=K_{1+((j-1)\bmod12)}\), qui désigne le
prolongement périodique également noté \(K_j\) dans
\eqref{eq:alpha-recurrencia}, le bilan antérieur au transport du
quotient est
\[
 Z_j:=P_j+E_j-\Phi_j-K_j^{\mathrm{per}}.
\]
La variable $s_j$ de \eqref{eq:alpha-recurrencia} désigne la somme qui
inclut déjà la contribution entrante, de sorte que
\[
 s_j=Z_j+c_{j+1},\qquad
 A_j=Z_j+c_{j+1}-1000c_j.
\]
La configuration incidencielle
$H^-_\alpha$ utilise les signes du bilan $Z_j$ de la fenêtre initiale ;
la détermination des chiffres utilise en outre les quotients de liaison
$c_{j+1}$ et $c_j$. Cette distinction conserve conjointement le support
signé et la normalisation du développement.
